% 第 1 章 绪论（以下为示例内容，撰写论文时请替换）
\chapter{绪论}

\section{研究背景与意义}

二阶椭圆型方程是偏微分方程理论中最基本的模型之一，广泛应用于稳态热传导、弹性力学与多孔介质渗流等问题。
由于解析解一般难以求得，数值方法成为研究这类方程的主要手段。有限元方法基于变分原理，
能够自然地处理复杂区域与边界条件，并具有较为完善的误差分析理论\cite{ciarlet1978,brenner2008}。

\section{国内外研究现状}

有限元方法的数学理论在二十世纪中后期逐步建立：变分形式与弱解理论可参见 \cite{evans2010,lax1954}，
Sobolev 空间工具可参见 \cite{adams2003}，而离散格式的一般分析框架可参见 \cite{strang1973}。
在此基础上，协调有限元、非协调有限元及自适应方法相继发展，误差估计与收敛性分析已成为标准工具。

\section{本文的主要工作}

本文围绕模型问题
\begin{equation}\label{eq:model}
-\Delta u + u = f \quad \text{in } \Omega, \qquad u = 0 \quad \text{on } \partial\Omega,
\end{equation}
完成以下工作：

\begin{enumerate}[label=(\arabic*)]
  \item 给出问题 \eqref{eq:model} 的弱形式，并证明弱解的存在唯一性；
  \item 构造分片线性协调有限元空间，推导能量范数与 $L^2$ 范数下的先验误差估计；
  \item 设计数值实验，验证收敛阶并讨论网格尺度对计算精度的影响。
\end{enumerate}

\section{论文结构安排}

第 2 章整理所需的函数空间与基本不等式；第 3 章给出主要定理及其证明；
第 4 章展示算法实现与数值结果；第 5 章总结全文并展望未来工作。
